% Propozycje podpisów do rozdz. 5 — NIE includowane w main.tex.
% Po akceptacji review: podmienić \caption w tex/5-badania.tex.

% --- Optuna / jakość (bez zmian modelu) ---
% fig:score_vs_hall
\caption{Zależność \emph{composite score} od wskaźnika halucynacji
         (\texttt{results-selected-9.05}). Flaga halucynacji determinuje
         score w~ponad 99\% przypadków.}

% fig:regression_r2
\caption{Porównanie $R^2$ modeli regresji wyjaśniających wariancję score
         (\texttt{results-selected-9.05}).}

% fig:nonhall_spread
\caption{Rozpiętość \emph{nonhall\_score} w~porównaniu z~\emph{composite score}
         (\texttt{results-selected-9.05}). Hiperparametry nie różnicują
         jakości odpowiedzi bez halucynacji.}

% fig:case_hall_rates
\caption{Wskaźnik halucynacji per przypadek testowy
         (\texttt{results-selected-9.05}).}

% fig:hparam_heatmap
\caption{Mapa korelacji hiperparametrów z~metrykami wynikowymi
         (\texttt{results-selected-9.05}). Brak silnych korelacji z~jakością.}

% fig:pareto_latency_hall  — Gauss mean±σ (nie mediana)
\caption{Kompromis Pareto: średnie opóźnienie E2E ($\pm\sigma$, model Gaussa)
         vs wskaźnik halucynacji (\texttt{results-selected-9.05}).
         Rozmiar punktu odpowiada średniemu \emph{composite score}.}

% fit: optuna_e2e_gauss_fit_matrix
\caption{Dopasowanie rozkładu normalnego do opóźnienia E2E
         (\texttt{results-selected-9.05}). Histogram oraz gęstość
         $\mathcal{N}(\hat\mu,\hat\sigma^2)$ per model.}

% fit: warm/cold generation LN
\caption{Dopasowanie rozkładu log-normalnego do \emph{generation\_ms}
         (warm/cold cache, $n\approx 100$). Histogram oraz gęstość LN
         per model i wariant kwantyzacji.}

% --- Warm kwantyzacja (log-normal) ---
% fig:dekompozycja_e2e
\caption{Dekompozycja E2E na TTFT i~\emph{generation\_ms}.
         Wysokości słupków to mediany modelu log-normalnego $e^{\hat\mu}$
         (warm cache, $n\approx 100$). Kwantyzacja wpływa na czas generacji,
         nie na prefill.}

% fig:speedup_q8
\caption{Przyspieszenie \emph{generation\_ms} względem Q8\_0
         w~eksperymencie warm-cache.
         Speedup $= e^{\hat\mu_{\mathrm{Q8}}}/e^{\hat\mu_{\mathrm{wariant}}}$
         (stosunek median modelu LN).}

% fig:gen_ms_variant
\caption{Czas \emph{generation\_ms} według wariantu kwantyzacji
         (mediana modelu LN $e^{\hat\mu}$, przedział $e^{\hat\mu\pm\hat\sigma}$;
         warm cache, $n\approx 100$).}

% fig:tok_s_variant
\caption{Szybkość generacji [znaki/s] (proxy tok/s) według wariantu kwantyzacji.
         Mediana modelu LN na per-rekord $\mathrm{answer\_chars}/(\mathrm{generation\_ms}/1000)$.}

% fig:gen_vs_length
\caption{Zależność czasu generacji od długości odpowiedzi
         (anomalia Q2\_K; warm cache, $n\approx 100$).}

% fig:ttft_warm
\caption{TTFT według wariantu kwantyzacji (warm cache):
         mediana modelu LN $e^{\hat\mu}$ z~przedziałem $e^{\hat\mu\pm\hat\sigma}$.}

% --- Cold (jeszcze nie w TeXu; do wstawienia później) ---
% cold_03_dekompozycja_e2e
\caption{Dekompozycja E2E na TTFT i~\emph{generation\_ms} (cold cache).
         Mediany modelu LN $e^{\hat\mu}$, $n\approx 100$.}

% cold_02_ttft / cold_07_generation / cold_04_speedup — analogicznie jak warm.
